% Results section for direct inclusion in a LaTeX paper
% Presents the five-step methodology outputs with cross-references
% to tables and figures generated by the implementation

\section{Results}\label{sec:results}

\subsection{Claim Extraction and Coverage (Steps 1--2)}

The claim extraction (Step~1) yields 41 claims from 56 clauses across
five applicable standards. Table~\ref{tab:standards} summarises the
standards and their roles in the integration.

The coverage matrix (Table~\ref{tab:coverage}) reveals asymmetric
coverage across lifecycle phases. The Concept/Requirements and
Verification phases are well-covered by multiple standards.
The Design/Training phase has limited coverage: ISO~26262 Parts~4 and~6
address conventional software design but are largely inapplicable to
neural network training (ISO~26262-6 Clause~8 prescribes coding guidelines
that PyTorch/TensorFlow frameworks violate structurally). Only ISO/PAS~8800
addresses AI-specific design. The Modification/Re-assurance phase has
the weakest coverage, with only partial contributions from three standards
and no complete workflow prescribed by any.

\subsection{Integrated GSN (Step 3)}

Table~\ref{tab:gsn-goals} presents the nine goals of the integrated
argument pattern. The extension from 6 to 9 goals is structural, not
cosmetic: G7 and G8 represent acceptability judgements about residual
risk (SOTIF and cybersecurity respectively) that are distinct from
verification activities. G9 is identified as undeveloped because no
standard prescribes a complete re-assurance workflow for AI model
modifications.

The goal density analysis (Table~\ref{tab:goal-density}) reveals that
G5 (V\&V sufficiency) is the only node where all four normative
standards contribute claims. G3 (data sufficiency) has the lowest density:
only ISO/PAS~8800 provides normative claims for training data quality,
with ISO/IEC~TR~5469 contributing informative guidance only.

The strategy S1 is reformulated from the base pattern to argue over four
assurance domains: (a)~AI-specific insufficiencies (ISO/PAS~8800),
(b)~functional insufficiencies (ISO~21448), (c)~cybersecurity risks
(ISO/SAE~21434), and (d)~hardware random/systematic faults (ISO~26262,
delegated via assumption A1.4).

\subsection{Requirement Inconsistencies (Step 4)}

Seven inconsistencies are identified at junction points in the integrated
GSN (Table~\ref{tab:inconsistencies}). Three are structural, two
terminological, and two methodological.

\paragraph{I-1: Incompatible risk classification frameworks.}
The top-level goal G1 inherits three risk classifications: ASIL~D from
ISO~26262 (severity~$\times$~exposure~$\times$~controllability),
qualitative acceptance criteria from ISO~21448 (project-specific), and
risk values 1--5 from ISO/SAE~21434 (impact~$\times$~attack feasibility).
These use different scales, different input parameters, and different
threshold definitions.

\paragraph{I-2: Evidence type asymmetry at G5 (central finding).}
Four evidence types converge at G5, each measuring a different property
on a different scale (Table~\ref{tab:evidence-types}). MC/DC coverage
(ISO~26262-6 Clause~9) is a binary structural metric. Scenario-based
testing (ISO~21448 Clauses~9--11) produces a scenario count with
triggering condition coverage. AI uncertainty quantification (ISO/PAS~8800
Clauses~8--9) produces a statistical distribution (e.g., AUROC for OOD
detection). Penetration testing (ISO/SAE~21434 Clause~10) produces
attack success rates. No standard defines how to combine these into a
single sufficiency claim.

\paragraph{I-5: Cybersecurity-SOTIF boundary.}
Adversarial point cloud perturbations cause false negatives that
simultaneously satisfy the definition of a SOTIF triggering condition
and a cybersecurity threat scenario. ISO~21448 explicitly excludes
cybersecurity threats (Clause~1). The boundary between G7 and G8 is
undefined.

\subsection{Assurance Gaps (Step 5)}

Six gaps are classified across four lifecycle phases
(Table~\ref{tab:gaps}). Three are missing evidence (ME), two are missing
claims (MC), and one is an unresolved inconsistency (UI).

\paragraph{Gap-1: No AI quantitative reliability target (ME).}
Hardware components at ASIL~D must meet SPFM~$\geq$~99\%,
LFM~$\geq$~90\%, and PMHF~$<$~$10^{-8}$~h$^{-1}$ (ISO~26262-5
Clause~9). No equivalent metric exists for AI model reliability.

\paragraph{Gap-2: No complete OTA re-assurance workflow (MC).}
G9 in the integrated GSN is undeveloped. ISO~26262-8 Clause~8 addresses
change management for conventional software. ISO/PAS~8800 Clause~14.8.3
provides partial re-approval guidance. Neither prescribes a complete
workflow for re-assuring an AI component after retraining.

\paragraph{Gap-3: Adversarial-SOTIF boundary unowned (UI, integration-induced).}
This gap is invisible when each standard is applied independently.
It appears only in the integrated argument where G7 (SOTIF) and G8
(cybersecurity) must coexist.

\paragraph{Gap-4: No cross-domain release criteria (MC, integration-induced).}
Four separate evidence sets (functional safety, SOTIF, cybersecurity,
AI-specific) must support a single release decision. No standard
prescribes combined evaluation criteria.

\subsection{CARLA Evaluation}

The evaluation under 25 weather combinations (Table~\ref{tab:weather-eval})
demonstrates the relationship between SOTIF triggering conditions and
detection performance. Under non-triggering conditions, mean recall
remains near the baseline level. Under triggering conditions, mean recall
decreases and geometric divergence increases, confirming that the deep
ensemble uncertainty metric correlates with environmental degradation.

The divergence metric provides evidence for two nodes in the integrated
GSN: G5 (V\&V evidence per ISO/PAS~8800) and G6 (monitoring evidence
per ISO/PAS~8800 Clause~14). This dual role demonstrates how a single
technical mechanism can address claims from multiple assurance domains.

The triggering condition coverage statistic (fraction of the weather
grid classified as triggering) provides evidence for ISO~21448
Clauses~9--11 at G5. In the case study, the SOTIF triggering fraction
is determined by the threshold definition (rain~$>$~20~mm/h or
visibility~$<$~200~m) applied to the parametric grid.
